\section{Results}

Every subsection below is labelled \textbf{[pre-declared]} where the analysis was specified before the results were seen, or \textbf{[post-hoc]} where it was added afterwards to explain something already observed. In summary: the protocol comparison (Section~4.1), the calibration analysis (Section~4.5) and the selective-prediction analysis (Section~4.6) were pre-declared; the surgeon-identifiability diagnostic (Section~4.2), the task-time comparison (Section~4.4) and the headline error metric itself (Section~4.3) were not. The distinction is load-bearing for a paper that examines protocol discipline, so it is carried in the text and not only in a table.

\subsection{The protocol gap [pre-declared]}

Skill classification accuracy under LOSO, where every surgeon appears in training, is 0.92 (suturing, 95\% interval 0.85-0.98), 0.97 (knot-tying, 0.92-1.00) and 0.93 (needle-passing, 0.85-1.00). Under LOUO, where the test surgeon is unseen, the same features and the same model give 0.44 (0.25-0.63), 0.31 (0.08-0.58) and 0.39 (0.16-0.69).

The drop is the paper's central observation, and it is best expressed as a paired difference. Bootstrapping the LOSO-minus-LOUO gap on the same resampled surgeons gives +0.49 (0.35-0.62) for suturing, +0.67 (0.41-0.89) for knot-tying and +0.54 (0.27-0.71) for needle-passing. All three intervals exclude zero by a wide margin. Changing only the evaluation protocol, with the features, the model and the data held constant, removes at least a third of the reported accuracy, and plausibly more than half; the point estimates are 0.49 to 0.67 but the intervals reach down to 0.27.

The LOUO point estimates fall at or below the rate obtainable by always predicting the most frequent class in that task --- 0.487, 0.444 and 0.393 respectively, so suturing and knot-tying fall below and needle-passing lands exactly on it. \textbf{We state this as a point estimate and not as a demonstration.} The LOUO intervals are wide enough to span the majority-class rate in all three tasks, so with eight surgeons we cannot establish that new-surgeon accuracy is below that floor; we can establish only that it is far below the LOSO figure. Quadratic-weighted kappa is negative under LOUO for suturing (-0.19) and needle-passing (-0.47), indicating agreement worse than chance on those two tasks, against +0.77 and +0.80 under LOSO.

% PROVENANCE: Figure 1A; Table 2. Source: kind=addendum_classification_ci, kind=addendum_protocol_gap,
\texttt{kind=addendum\_majority\_class}, \texttt (Table T2).*

\subsection{The same features identify the surgeon [post-hoc]}

This analysis was added after observing the gap in Section~4.1, in order to explain it. It was not part of the pre-registered plan, and it should be read as a diagnostic rather than as a confirmatory test.

Substituting surgeon identity for skill as the target, and leaving the features, the model, the folds and the regularisation unchanged, the same 16 features identify which of eight surgeons performed a held-out repetition with accuracy 0.87 (95\% interval 0.77-0.97) for suturing, 0.83 (0.67-0.97) for knot-tying and 0.93 (0.83-1.00) for needle-passing, against chance rates of 0.125, 0.125 and 0.143. All three intervals lie far above chance.

Because the skill label is a property of the surgeon and is constant across that surgeon's trials, a model that can identify the surgeon can retrieve the label without representing skill. This is a sufficient mechanism for the LOSO figures in Section~4.1, and it is consistent with their collapse under LOUO, where identity carries no usable information about an unseen individual. It is not proof that the classifier relies on identity: we have demonstrated that the information is present and recoverable, not that this particular model uses it. Distinguishing the two would require an intervention on the features, which we did not perform.

% PROVENANCE: Figure 1B; Table 2. Source: kind=addendum_surgeon_identity_ci,
\texttt{kind=diagnostic\_surgeon\_identity}.*

\subsection{GRS prediction against a feature-free reference [pre-declared design, post-hoc metric]}

The design of this analysis was pre-registered; the metric reported was not, for the reason given in Section~3.7. We flag that explicitly because the substitution is favourable to readability and could otherwise pass unnoticed.

Under LOUO, ridge regression reduces mean absolute error relative to the feature-free reference by +1.54 GRS points (95\% interval 0.65-2.52) on knot-tying, from 5.10 to 3.56. On suturing the reduction is +0.15 (-1.04 to 1.34) and on needle-passing it is -0.50 (-1.21 to -0.04) --- that is, on needle-passing the fitted model is \textit{worse} than predicting the training mean, with an interval excluding zero on the wrong side. Under the pre-declared rule, only knot-tying beats the reference. The zero-shot path behaves the same way: +2.30 (1.14-3.45) on knot-tying, -0.05 (-1.05 to 0.96) on suturing and -0.24 (-0.79 to 0.31) on needle-passing. The two knot-tying reductions are not distinguishable from one another: their intervals overlap across almost their whole length, and nothing in this study ranks the two paths. Its fitted slope was positive in all 38 folds, so the direction guard of Section~3.5 passed and no inverted score is being rescued by the calibration line.

Under LOSO the same comparison is positive on all three tasks for ridge (+1.15, +1.78, +0.43, all intervals excluding zero), which is the Section~4.1 pattern reappearing in the regression target.

Agreement statistics tell a consistent story. Under LOUO, ICC(2,1) is 0.60 (0.35-0.78) for ridge and 0.73 (0.53-0.85) for the zero-shot path on knot-tying, and is indistinguishable from zero on the other two tasks for both paths (ridge 0.20 and 0.00; zero-shot 0.06 and -0.13). Bland-Altman bias is small everywhere (-0.37 to +0.52 GRS points), but the limits of agreement span roughly ±11 points on suturing and needle-passing against ±7 on knot-tying --- that is, on two of three tasks the disagreement between model and rater is about as large as the full spread of the scale. For context, pooled inter-rater reliability for global rating scales in simulation-based assessment has been reported at 0.78 (95\% CI 0.71-0.83, 23 studies; Ilgen et al. 2015), with the original OSATS report giving 0.64-0.72 \cite{martin1997osats}. Neither source names the coefficient it pooled, so neither is strictly an ICC(2,1) and both comparisons are approximate. Our strongest agreement, 0.73 for the zero-shot path on knot-tying, sits below the pooled figure and carries an interval (0.53-0.85) wide enough to be uninformative about the comparison. We avoid the language of a ceiling: rater reliability is not a fixed bar but depends heavily on training, with untrained raters reported at ICC 0.26, rising to 0.76 immediately after training but varying widely on later re-testing (-0.17 to 0.70), which the authors summarise as the benefit of training being ``lost in periods of nonuse'' \cite{gawad2019interrater}. The target these models are fitted to is therefore itself noisy, and the reading we take is simply that no method here demonstrated agreement comparable to trained human raters on any task.

We note that the pooled out-of-fold Spearman correlations for these same predictions (LOUO: 0.21, 0.64, -0.13 for ridge) must be read against the feature-free reference's values on the identical folds (-0.78, -0.86, -0.79) rather than against zero, for the reason established in Section~3.7.

% PROVENANCE: Figure 1C; Table 3. Source: kind=headline_mae_vs_reference, kind=jev_headline_mae_vs_reference,
\texttt{kind=eval\_agreement}, \texttt *

\subsection{A one-feature comparison [post-hoc]}

This comparison was added after Section~4.3, to ask what the surviving knot-tying signal consists of. It is diagnostic.

The zero-shot model's overall score correlates with total task time at Spearman -0.91 (suturing), -0.81 (knot-tying) and -0.86 (needle-passing). Correlations with all 16 features are logged, so this is not a selected maximum: task time is the strongest associate in every task, with the next strongest between 0.41 and 0.71 in absolute value.

Substituting a single ordinary-least-squares line on task time alone for either model, on identical folds and under the identical metric, reduces error by +1.88 GRS points (0.94-2.89) on knot-tying under LOUO. On suturing (+0.10, -0.72 to 0.83) and needle-passing (-0.29, -0.88 to 0.11) it fails to beat the reference, as both models do. The three knot-tying intervals --- 0.65-2.52 for ridge, 1.14-3.45 for the zero-shot path and 0.94-2.89 for task time alone --- overlap across almost their whole length. We therefore report them as three estimates of one effect and draw no ordering among them; the point estimates differ, but this study cannot resolve that difference.

On this evidence the surviving cross-surgeon signal is, to within the resolution of eight surgeons, the signal available from trial duration alone. We do not claim the two are equivalent; we claim that this study does not demonstrate any advantage for either the 16-feature model or the zero-shot model over a one-feature baseline on the one task where a signal survives.

% PROVENANCE: Source: kind=addendum_jev_score_vs_feature, kind=posthoc_time_only_headline.
\subsection{Calibration [pre-declared]}

Each path is assessed against the target it was given: the baseline classifier against the self-declared skill label it was trained on, and the zero-shot model against the level implied by GRS, raw and with no recalibration applied. The two are therefore not compared with one another.

On a held-out surgeon, the baseline classifier states a mean confidence of 0.755 while achieving 0.379 accuracy. Its expected calibration error is 0.377 (0.237-0.541), and its Brier score of 0.875 is worse than the 0.647 obtained by a constant forecast equal to the observed class frequencies --- a Brier skill score of -0.351, with an interval (-0.688 to -0.092) lying entirely below zero. Its most confident bin, containing 46 of 103 trials at a mean stated confidence of 0.92, is correct 39\% of the time (\texttt{addendum\_reliability\_bin}).

Under LOSO the same classifier appears well calibrated and, if anything, underconfident: mean confidence 0.837 against accuracy 0.942, expected calibration error 0.105, Brier skill +0.778 (0.698-0.856). The same model and the same features thus yield opposite calibration verdicts under the two protocols, which is Section~4.1 restated in the confidence domain.

The zero-shot model is also overconfident: mean confidence 0.696 against accuracy 0.301, expected calibration error 0.395 (0.317-0.494), Brier skill -0.433 (-0.588 to -0.242). Its predictions are usually close --- within one level on 84\% of trials --- but they sit systematically high: the mean predicted level is 2.05 against a mean true level of 1.71, an offset of about a third of a level, and it places 39 trials in the top two levels where the GRS-implied truth places 22 (\texttt{addendum\_jev\_level\_agreement}, \texttt{addendum\_jev\_level\_usage}). That offset means its exact-level miscalibration mixes genuine overconfidence with a scale mismatch between its rubric and our GRS-derived levels; the two cannot be separated without a recalibration that was excluded in advance, and we therefore do not attribute the full -0.433 to overconfidence.

% PROVENANCE: Figure 2 (eval_reliability.png); Table 4. Source: kind=eval_calibration,
\texttt{kind=addendum\_mean\_confidence}, \texttt{kind=addendum\_reliability\_bin}.*

\subsection{Selective prediction [pre-declared]}

The deployment question is whether a model's confidence identifies the trials on which it should defer to a human rater. We rank the 103 pooled trials by each path's own confidence, defer the least confident first, and measure the absolute GRS error of those retained.

\textbf{The two paths do not support the same inference here, and we state that before the numbers.} Ridge regression emits no native confidence, so some proxy was required; we declared in advance that the baseline would pair its regression prediction with the classifier's top-class probability. Those are two different models fitted to two different targets, and Section~3.1 establishes that the targets disagree --- one surgeon labelled expert averages GRS 13.5 while one labelled novice averages 19.1. The baseline's confidence therefore carries no information about the regression's error except whatever the two targets happen to share. \textbf{Its null result is uninformative about whether ridge knows when it is wrong; it is a property of the pairing.} The zero-shot path has no such problem: its score and its confidence come from one model answering one question, so it is the only clean test of self-knowledge in this study.

The clean test is null. Ordering by the zero-shot model's own confidence lowers the area under the risk-coverage curve by 0.155 GRS points relative to random ordering under LOUO, with a 95\% interval of -0.416 to 0.606; substituting its maximum level probability for its reported confidence changes nothing (0.149). Under LOSO the figure is 0.077. The uninterpretable arm is also null, for the record: 0.135 (-0.628 to 0.922) under LOUO and 0.045 under LOSO. Expressed as a fraction of the distance between random and ideal ordering these are 7.4\% and 5.5\% under LOUO, and 3.9\% and 2.5\% under LOSO --- percentages of a quantity that is itself indistinguishable from zero, so they should be read only alongside the intervals above and not quoted alone.

The coverage table shows the same flatness directly. Deferring 20\% of trials by baseline confidence under LOUO changes the retained error from 4.31 to 4.50 GRS points --- it rises --- and deferring 60\% gives 4.15. The zero-shot path falls from 4.05 to 3.76 at 40\% coverage, a change no interval supports.

One further caution. The upper interval bounds are +0.61 GRS points for the clean arm and +0.92 for the other, up to a fifth of the overall error, so this is an under-powered null rather than a demonstration that confidence is useless. In summary, this study yields one clean null and one uninterpretable one, not two independent confirmations. A predictor with native uncertainty --- quantile or Bayesian regression, or a conformal wrapper --- would remove the mismatch and is the concrete design a follow-up should adopt.

% PROVENANCE: Figure 3 (eval_risk_coverage.png); Table 5. Source: kind=eval_selective, kind=eval_coverage.
\subsection{Robustness}

Conclusions do not depend on the filtering choice. Repeating the LOUO analysis on unfiltered features and at an 8 Hz cutoff leaves knot-tying positive throughout (+1.18 to +1.69 GRS points) and suturing and needle-passing non-significant throughout; knot-tying's interval excludes zero in four of the six configuration-by-filter combinations, failing at 8 Hz (lower bound -0.26). Widening the ridge penalty grid in a second configuration leaves knot-tying unchanged and moves suturing and needle-passing toward the reference, because the inner search then selects maximal shrinkage; both configurations are reported.

The zero-shot model is close to, but not exactly, deterministic: three repetitions of an identical request agree at Spearman 0.999 with a mean absolute score difference of 0.022-0.024 on a 0-4 scale, and only 19-21\% of trials return a bit-identical score (\texttt{addendum\_jev\_repeatability}). Averaging the three runs changes every headline figure by at most 0.01 GRS points.

% PROVENANCE: Source: kind=headline_mae_vs_reference and kind=jev_headline_mae_vs_reference across
\texttt{feature\_set} values \texttt{smoothed\_5hz}, \texttt{raw} and \texttt{smoothed\_8hz}; \texttt across both configuration files.*

---
